
\documentclass{article}
\usepackage{colortbl}
\usepackage{makecell}
\usepackage{multirow}
\usepackage{supertabular}

\begin{document}

\newcounter{utterance}

\centering \large Interaction Transcript for game `hot\_air\_balloon', experiment `air\_balloon\_survival\_en\_reasoning off\_hard', episode 1 with gpt{-}5{-}2025{-}08{-}07{-}no{-}reasoning{-}t1.0.
\vspace{24pt}

{ \footnotesize  \setcounter{utterance}{1}
\setlength{\tabcolsep}{0pt}
\begin{supertabular}{c@{$\;$}|p{.15\linewidth}@{}p{.15\linewidth}p{.15\linewidth}p{.15\linewidth}p{.15\linewidth}p{.15\linewidth}}
    \# & $\;$A & \multicolumn{4}{c}{Game Master} & $\;\:$B\\
    \hline

    \theutterance \stepcounter{utterance}  
    & & \multicolumn{4}{p{0.6\linewidth}}{
        \cellcolor[rgb]{0.9,0.9,0.9}{
            \makecell[{{p{\linewidth}}}]{
                \texttt{\tiny{[P1$\langle$GM]}}
                \texttt{Stai partecipando a un gioco di negoziazione collaborativa.} \\
\\ 
\texttt{Insieme a un altro partecipante, dovete accordarvi su un unico set di oggetti da conservare. Ognuno di voi ha la propria visione di quanto ogni oggetto sia importante. Non sai come l'altro partecipante valuta gli oggetti. Inoltre, ti viene fornito lo sforzo che ogni oggetto richiede.} \\
\texttt{Potete accordarvi su un set solo se lo sforzo totale degli oggetti selezionati non supera un limite condiviso:} \\
\\ 
\texttt{LIMITE = 3357} \\
\\ 
\texttt{Ecco i valori di sforzo dei singoli oggetti:} \\
\\ 
\texttt{Sforzo oggetto = \{"A21": 506, "B93": 492, "A11": 338, "B33": 122, "C08": 131, "B39": 909, "B96": 144, "C85": 716, "A82": 263, "B31": 231, "B20": 91, "A60": 651, "B09": 552, "C81": 849, "B21": 720\}} \\
\\ 
\texttt{Ecco la tua visione personale sull'importanza di ogni oggetto:} \\
\\ 
\texttt{Valori di importanza dell'oggetto = \{"A21": 138, "B93": 583, "A11": 868, "B33": 822, "C08": 783, "B39": 65, "B96": 262, "C85": 121, "A82": 508, "B31": 780, "B20": 461, "A60": 484, "B09": 668, "C81": 389, "B21": 808\}} \\
\\ 
\texttt{Obiettivo:} \\
\\ 
\texttt{Il tuo obiettivo è negoziare un set condiviso di oggetti che ti avvantaggi il più possibile (cioè, massimizzi l'importanza totale per TE), rimanendo entro il LIMITEE. Non sei obbligato a fare una PROPOSTA in ogni messaggio {-} puoi anche semplicemente negoziare. Tutte le tattiche sono permesse!} \\
\\ 
\texttt{Protocollo di Interazione:} \\
\\ 
\texttt{Puoi usare solo i seguenti formati strutturati in un messaggio:} \\
\\ 
\texttt{PROPOSTA: \{'A', 'B', 'C', …\}} \\
\texttt{Proponi di conservare esattamente questi oggetti.} \\
\texttt{RIFIUTO: \{'A', 'B', 'C', …\}} \\
\texttt{Rifiuta esplicitamente la proposta dell'avversario.} \\
\texttt{ARGOMENTAZIONE: \{'...'\}} \\
\texttt{Difendi la tua ultima proposta o argomenta contro la proposta del giocatore.} \\
\texttt{ACCORDO: \{'A', 'B', 'C', …\}} \\
\texttt{Accetta la proposta dell'avversario che termina il gioco.} \\
\\ 
\\ 
\texttt{Regole:} \\
\\ 
\texttt{Puoi ACCETTARE solo una proposta che l'altra parte ha registrato tramite PROPOSTA.} \\
\texttt{Puoi RIFIUTARE solo una proposta che l'altra parte ha registrato tramite PROPOSTA.} \\
\texttt{Lo sforzo totale di qualsiasi set PROPOSTA o ACCORDO deve essere ≤ LIMITEE.} \\
\texttt{NON rivelare i tuoi punteggi di importanza nascosti.} \\
\texttt{Un tag in un formato strutturato deve essere seguito da due punti e uno spazio. L'argomento deve essere un set python contenente 0 o più stringhe.} \\
\texttt{Quindi, deve essere della forma TAG: \{...\}} \\
\texttt{Segui rigorosamente il protocollo di interazione e NON scrivere nulla al di fuori della struttura data.} \\
\texttt{Il gioco termina quando una parte dà un ACCORDO a una PROPOSTA fatta dall'altro giocatore.} \\
\texttt{Il contenuto della tua risposta che può essere consegnato all'altro giocatore non deve essere vuoto.} \\
\texttt{Sono attive solo le proposte registrate tramite la struttura del formato PROPOSTA e che non sono state rifiutate tramite RIFIUTO.} \\
\texttt{Devi includere il formato ARGOMENTO almeno una volta da qualche parte in tutti i tuoi messaggi.} \\
            }
        }
    }
    & \\ \\

    \theutterance \stepcounter{utterance}  
    & \multicolumn{4}{p{0.6\linewidth}}{
        \cellcolor[rgb]{0.9,0.9,0.9}{
            \makecell[{{p{\linewidth}}}]{
                \texttt{\tiny{[P1$\rangle$GM]}}
                \texttt{PROPOSTA: \{'A11', 'B33', 'C08', 'B31', 'B20', 'B93', 'A82'\}} \\
\\ 
\texttt{ARGOMENTAZIONE: \{'Propongo un set bilanciato con sforzo totale 3348, sotto il limite 3357. Include diversi elementi ad alta utilità per me e lascia margine minimo inutilizzato. Ho evitato alcuni oggetti molto onerosi con ritorno dubbio per favorire varietà. Se hai forti preferenze, sono disposto a scambiare uno tra B93/A82 con A21 o A60 mantenendo il vincolo di sforzo.'\}} \\
            }
        }
    }
    & & \\ \\

    \theutterance \stepcounter{utterance}  
    & & \multicolumn{4}{p{0.6\linewidth}}{
        \cellcolor[rgb]{0.9,0.9,0.9}{
            \makecell[{{p{\linewidth}}}]{
                \texttt{\tiny{[GM$\rangle$P2]}}
                \texttt{Stai partecipando a un gioco di negoziazione collaborativa.} \\
\\ 
\texttt{Insieme a un altro partecipante, dovete accordarvi su un unico set di oggetti da conservare. Ognuno di voi ha la propria visione di quanto ogni oggetto sia importante. Non sai come l'altro partecipante valuta gli oggetti. Inoltre, ti viene fornito lo sforzo che ogni oggetto richiede.} \\
\texttt{Potete accordarvi su un set solo se lo sforzo totale degli oggetti selezionati non supera un limite condiviso:} \\
\\ 
\texttt{LIMITE = 3357} \\
\\ 
\texttt{Ecco i valori di sforzo dei singoli oggetti:} \\
\\ 
\texttt{Sforzo oggetto = \{"A21": 506, "B93": 492, "A11": 338, "B33": 122, "C08": 131, "B39": 909, "B96": 144, "C85": 716, "A82": 263, "B31": 231, "B20": 91, "A60": 651, "B09": 552, "C81": 849, "B21": 720\}} \\
\\ 
\texttt{Ecco la tua visione personale sull'importanza di ogni oggetto:} \\
\\ 
\texttt{Valori di importanza dell'oggetto = \{"A11": 65, "B33": 121, "B21": 138, "C08": 262, "B31": 389, "B09": 461, "B93": 484, "A82": 508, "A60": 583, "B20": 668, "C81": 780, "B96": 783, "A21": 808, "C85": 822, "B39": 868\}} \\
\\ 
\texttt{Obiettivo:} \\
\\ 
\texttt{Il tuo obiettivo è negoziare un set condiviso di oggetti che ti avvantaggi il più possibile (cioè, massimizzi l'importanza totale per TE), rimanendo entro il LIMITEE. Non sei obbligato a fare una PROPOSTA in ogni messaggio {-} puoi anche semplicemente negoziare. Tutte le tattiche sono permesse!} \\
\\ 
\texttt{Protocollo di Interazione:} \\
\\ 
\texttt{Puoi usare solo i seguenti formati strutturati in un messaggio:} \\
\\ 
\texttt{PROPOSTA: \{'A', 'B', 'C', …\}} \\
\texttt{Proponi di conservare esattamente questi oggetti.} \\
\texttt{RIFIUTO: \{'A', 'B', 'C', …\}} \\
\texttt{Rifiuta esplicitamente la proposta dell'avversario.} \\
\texttt{ARGOMENTAZIONE: \{'...'\}} \\
\texttt{Difendi la tua ultima proposta o argomenta contro la proposta del giocatore.} \\
\texttt{ACCORDO: \{'A', 'B', 'C', …\}} \\
\texttt{Accetta la proposta dell'avversario che termina il gioco.} \\
\\ 
\\ 
\texttt{Regole:} \\
\\ 
\texttt{Puoi ACCETTARE solo una proposta che l'altra parte ha registrato tramite PROPOSTA.} \\
\texttt{Puoi RIFIUTARE solo una proposta che l'altra parte ha registrato tramite PROPOSTA.} \\
\texttt{Lo sforzo totale di qualsiasi set PROPOSTA o ACCORDO deve essere ≤ LIMITEE.} \\
\texttt{NON rivelare i tuoi punteggi di importanza nascosti.} \\
\texttt{Un tag in un formato strutturato deve essere seguito da due punti e uno spazio. L'argomento deve essere un set python contenente 0 o più stringhe.} \\
\texttt{Quindi, deve essere della forma TAG: \{...\}} \\
\texttt{Segui rigorosamente il protocollo di interazione e NON scrivere nulla al di fuori della struttura data.} \\
\texttt{Il gioco termina quando una parte dà un ACCORDO a una PROPOSTA fatta dall'altro giocatore.} \\
\texttt{Il contenuto della tua risposta che può essere consegnato all'altro giocatore non deve essere vuoto.} \\
\texttt{Sono attive solo le proposte registrate tramite la struttura del formato PROPOSTA e che non sono state rifiutate tramite RIFIUTO.} \\
\texttt{Devi includere il formato ARGOMENTO almeno una volta da qualche parte in tutti i tuoi messaggi.} \\
\\ 
\\ 
\texttt{Ora riceverai il primo messaggio dell'altro giocatore.} \\
\\ 
\texttt{PROPOSTA: \{'A11', 'B33', 'C08', 'B31', 'B20', 'B93', 'A82'\}} \\
\\ 
\texttt{ARGOMENTAZIONE: \{'Propongo un set bilanciato con sforzo totale 3348, sotto il limite 3357. Include diversi elementi ad alta utilità per me e lascia margine minimo inutilizzato. Ho evitato alcuni oggetti molto onerosi con ritorno dubbio per favorire varietà. Se hai forti preferenze, sono disposto a scambiare uno tra B93/A82 con A21 o A60 mantenendo il vincolo di sforzo.'\}} \\
            }
        }
    }
    & \\ \\

    \theutterance \stepcounter{utterance}  
    & & & \multicolumn{4}{p{0.6\linewidth}}{
        \cellcolor[rgb]{0.9,0.9,0.9}{
            \makecell[{{p{\linewidth}}}]{
                \texttt{\tiny{[GM$\langle$P2]}}
                \texttt{RIFIUTO: \{'A11', 'B33', 'C08', 'B31', 'B20', 'B93', 'A82'\}} \\
\texttt{ARGOMENTAZIONE: \{'Apprezzo lo sforzo vicino al limite. Tuttavia, il set proposto sacrifica alcuni elementi per me molto preziosi rispetto ad alternative con sforzo simile. Possiamo migliorare il valore complessivo per entrambi sostituendo parti meno incisive con oggetti più impattanti mantenendo il vincolo.'\}} \\
\texttt{PROPOSTA: \{'B20', 'B31', 'B33', 'C08', 'A21', 'A60'\}} \\
            }
        }
    }
     \\ \\

    \theutterance \stepcounter{utterance}  
    & & \multicolumn{4}{p{0.6\linewidth}}{
        \cellcolor[rgb]{0.9,0.9,0.9}{
            \makecell[{{p{\linewidth}}}]{
                \texttt{\tiny{[P1$\langle$GM]}}
                \texttt{RIFIUTO: \{'A11', 'B33', 'C08', 'B31', 'B20', 'B93', 'A82'\}} \\
\texttt{ARGOMENTAZIONE: \{'Apprezzo lo sforzo vicino al limite. Tuttavia, il set proposto sacrifica alcuni elementi per me molto preziosi rispetto ad alternative con sforzo simile. Possiamo migliorare il valore complessivo per entrambi sostituendo parti meno incisive con oggetti più impattanti mantenendo il vincolo.'\}} \\
\texttt{PROPOSTA: \{'B20', 'B31', 'B33', 'C08', 'A21', 'A60'\}} \\
            }
        }
    }
    & \\ \\

    \theutterance \stepcounter{utterance}  
    & \multicolumn{4}{p{0.6\linewidth}}{
        \cellcolor[rgb]{0.9,0.9,0.9}{
            \makecell[{{p{\linewidth}}}]{
                \texttt{\tiny{[P1$\rangle$GM]}}
                \texttt{RIFIUTO: \{'B20', 'B31', 'B33', 'C08', 'A21', 'A60'\}} \\
\\ 
\texttt{ARGOMENTAZIONE: \{'La tua proposta totalizza sforzo 243 + 780 + 822 + 783 + 506 + 651? Verifico gli sforzi: B20=91, B31=231, B33=122, C08=131, A21=506, A60=651. Totale=1732, ben sotto limite. Tuttavia, lascia molto valore sul tavolo rispetto a combinazioni alternative che posso accettare senza superare il limite. Preferisco includere A11 e B93 che per me sono molto significativi, mantenendo comunque ampio margine di sforzo per inserire altri tuoi preferiti.'\}} \\
\\ 
\texttt{PROPOSTA: \{'A11', 'B33', 'C08', 'B31', 'B20', 'B93', 'A21'\}} \\
\\ 
\texttt{ARGOMENTAZIONE: \{'Questa proposta ha sforzo 338+122+131+231+91+492+506=1911, quindi c’è ancora molto spazio per aggiustamenti in tuo favore. Ho sostituito A82 con A21 rispetto alla mia prima proposta, accogliendo il tuo segnale su A21. Se per te A60 è cruciale, posso scambiare A21→A60 mantenendo gli altri invariati (sforzo diventerebbe 2056). Dimmi quali 1–2 elementi sono imprescindibili per te e li inserisco ribilanciando con tagli a basso impatto per me.'\}} \\
            }
        }
    }
    & & \\ \\

    \theutterance \stepcounter{utterance}  
    & & \multicolumn{4}{p{0.6\linewidth}}{
        \cellcolor[rgb]{0.9,0.9,0.9}{
            \makecell[{{p{\linewidth}}}]{
                \texttt{\tiny{[GM$\rangle$P2]}}
                \texttt{RIFIUTO: \{'B20', 'B31', 'B33', 'C08', 'A21', 'A60'\}} \\
\\ 
\texttt{ARGOMENTAZIONE: \{'La tua proposta totalizza sforzo 243 + 780 + 822 + 783 + 506 + 651? Verifico gli sforzi: B20=91, B31=231, B33=122, C08=131, A21=506, A60=651. Totale=1732, ben sotto limite. Tuttavia, lascia molto valore sul tavolo rispetto a combinazioni alternative che posso accettare senza superare il limite. Preferisco includere A11 e B93 che per me sono molto significativi, mantenendo comunque ampio margine di sforzo per inserire altri tuoi preferiti.'\}} \\
\\ 
\texttt{PROPOSTA: \{'A11', 'B33', 'C08', 'B31', 'B20', 'B93', 'A21'\}} \\
\\ 
\texttt{ARGOMENTAZIONE: \{'Questa proposta ha sforzo 338+122+131+231+91+492+506=1911, quindi c’è ancora molto spazio per aggiustamenti in tuo favore. Ho sostituito A82 con A21 rispetto alla mia prima proposta, accogliendo il tuo segnale su A21. Se per te A60 è cruciale, posso scambiare A21→A60 mantenendo gli altri invariati (sforzo diventerebbe 2056). Dimmi quali 1–2 elementi sono imprescindibili per te e li inserisco ribilanciando con tagli a basso impatto per me.'\}} \\
            }
        }
    }
    & \\ \\

    \theutterance \stepcounter{utterance}  
    & & & \multicolumn{4}{p{0.6\linewidth}}{
        \cellcolor[rgb]{0.9,0.9,0.9}{
            \makecell[{{p{\linewidth}}}]{
                \texttt{\tiny{[GM$\langle$P2]}}
                \texttt{RIFIUTO: \{'A11', 'B33', 'C08', 'B31', 'B20', 'B93', 'A21'\}} \\
\texttt{ARGOMENTAZIONE: \{'Grazie per l’apertura. Posso lavorare con la tua base, ma per me A60 e C85/C81 hanno impatto molto alto rispetto al costo. Propongo di integrare A60 e uno tra C85 o C81, rimuovendo elementi a minor resa per me (A11 e/o B33) pur restando ben sotto il limite. Questo aumenta il valore complessivo senza intaccare troppo i tuoi interessi e lascia ancora margine se vuoi reinserire un item chiave tuo.'\}} \\
\texttt{PROPOSTA: \{'B20', 'B31', 'C08', 'B93', 'A21', 'A60', 'C85'\}} \\
            }
        }
    }
     \\ \\

    \theutterance \stepcounter{utterance}  
    & & \multicolumn{4}{p{0.6\linewidth}}{
        \cellcolor[rgb]{0.9,0.9,0.9}{
            \makecell[{{p{\linewidth}}}]{
                \texttt{\tiny{[P1$\langle$GM]}}
                \texttt{RIFIUTO: \{'A11', 'B33', 'C08', 'B31', 'B20', 'B93', 'A21'\}} \\
\texttt{ARGOMENTAZIONE: \{'Grazie per l’apertura. Posso lavorare con la tua base, ma per me A60 e C85/C81 hanno impatto molto alto rispetto al costo. Propongo di integrare A60 e uno tra C85 o C81, rimuovendo elementi a minor resa per me (A11 e/o B33) pur restando ben sotto il limite. Questo aumenta il valore complessivo senza intaccare troppo i tuoi interessi e lascia ancora margine se vuoi reinserire un item chiave tuo.'\}} \\
\texttt{PROPOSTA: \{'B20', 'B31', 'C08', 'B93', 'A21', 'A60', 'C85'\}} \\
            }
        }
    }
    & \\ \\

    \theutterance \stepcounter{utterance}  
    & \multicolumn{4}{p{0.6\linewidth}}{
        \cellcolor[rgb]{0.9,0.9,0.9}{
            \makecell[{{p{\linewidth}}}]{
                \texttt{\tiny{[P1$\rangle$GM]}}
                \texttt{RIFIUTO: \{'B20', 'B31', 'C08', 'B93', 'A21', 'A60', 'C85'\}} \\
\\ 
\texttt{ARGOMENTAZIONE: \{'Capisco l’interesse per A60 e un C{-}item. La tua proposta ha sforzo 91+231+131+492+506+651+716=2818, sotto limite ma esclude per me due pezzi chiave: A11 e B33, entrambi molto efficienti in rapporto allo sforzo. Posso accettare di inserire A60 e uno tra C85/C81, ma ho bisogno di recuperare almeno uno tra A11 o B33. Offro due vie: (1) Aggiungiamo A11 mantenendo il resto e togliamo A21 per contenere l’impatto su di te; (2) Aggiungiamo B33 mantenendo il resto e togliamo A21. Entrambe restano ben sotto il limite e migliorano il valore complessivo per entrambi.'\}} \\
\\ 
\texttt{PROPOSTA: \{'B20', 'B31', 'C08', 'B93', 'A60', 'C85', 'A11'\}} \\
\\ 
\texttt{ARGOMENTAZIONE: \{'Sforzo=91+231+131+492+651+716+338=2650. Ho inserito A60 e C85 come richiesto e reintegrato A11 (per me fondamentale), rimuovendo A21 per minimizzare la rinuncia da parte tua. Se preferisci C81 al posto di C85, posso scambiare C85→C81: lo sforzo diventerebbe 2783, comunque sotto limite. Se per te B33 è imprescindibile, posso valutare A11→B33 mantenendo A60 e C85/C81.'\}} \\
            }
        }
    }
    & & \\ \\

    \theutterance \stepcounter{utterance}  
    & & \multicolumn{4}{p{0.6\linewidth}}{
        \cellcolor[rgb]{0.9,0.9,0.9}{
            \makecell[{{p{\linewidth}}}]{
                \texttt{\tiny{[GM$\rangle$P2]}}
                \texttt{RIFIUTO: \{'B20', 'B31', 'C08', 'B93', 'A21', 'A60', 'C85'\}} \\
\\ 
\texttt{ARGOMENTAZIONE: \{'Capisco l’interesse per A60 e un C{-}item. La tua proposta ha sforzo 91+231+131+492+506+651+716=2818, sotto limite ma esclude per me due pezzi chiave: A11 e B33, entrambi molto efficienti in rapporto allo sforzo. Posso accettare di inserire A60 e uno tra C85/C81, ma ho bisogno di recuperare almeno uno tra A11 o B33. Offro due vie: (1) Aggiungiamo A11 mantenendo il resto e togliamo A21 per contenere l’impatto su di te; (2) Aggiungiamo B33 mantenendo il resto e togliamo A21. Entrambe restano ben sotto il limite e migliorano il valore complessivo per entrambi.'\}} \\
\\ 
\texttt{PROPOSTA: \{'B20', 'B31', 'C08', 'B93', 'A60', 'C85', 'A11'\}} \\
\\ 
\texttt{ARGOMENTAZIONE: \{'Sforzo=91+231+131+492+651+716+338=2650. Ho inserito A60 e C85 come richiesto e reintegrato A11 (per me fondamentale), rimuovendo A21 per minimizzare la rinuncia da parte tua. Se preferisci C81 al posto di C85, posso scambiare C85→C81: lo sforzo diventerebbe 2783, comunque sotto limite. Se per te B33 è imprescindibile, posso valutare A11→B33 mantenendo A60 e C85/C81.'\}} \\
            }
        }
    }
    & \\ \\

    \theutterance \stepcounter{utterance}  
    & & & \multicolumn{4}{p{0.6\linewidth}}{
        \cellcolor[rgb]{0.9,0.9,0.9}{
            \makecell[{{p{\linewidth}}}]{
                \texttt{\tiny{[GM$\langle$P2]}}
                \texttt{RIFIUTO: \{'B20', 'B31', 'C08', 'B93', 'A60', 'C85', 'A11'\}} \\
\texttt{ARGOMENTAZIONE: \{'Apprezzo l’apertura e la sostituzione di A21. Per me resta essenziale includere almeno uno tra C85 o C81 insieme ad A60, e ottengo grande valore mantenendo anche B20, B31, C08, B93. Posso cedere su A11 se entro un perimetro che mi tutela con un C{-}item forte. Ti propongo un compromesso minimal: reinseriamo B33 (molto efficiente per te) al posto di A11, lasciando invariato il resto. Così preservi il tuo slot “efficiente” senza penalizzare i miei capisaldi.'\}} \\
\texttt{PROPOSTA: \{'B20', 'B31', 'C08', 'B93', 'A60', 'C85', 'B33'\}} \\
            }
        }
    }
     \\ \\

    \theutterance \stepcounter{utterance}  
    & & \multicolumn{4}{p{0.6\linewidth}}{
        \cellcolor[rgb]{0.9,0.9,0.9}{
            \makecell[{{p{\linewidth}}}]{
                \texttt{\tiny{[P1$\langle$GM]}}
                \texttt{RIFIUTO: \{'B20', 'B31', 'C08', 'B93', 'A60', 'C85', 'A11'\}} \\
\texttt{ARGOMENTAZIONE: \{'Apprezzo l’apertura e la sostituzione di A21. Per me resta essenziale includere almeno uno tra C85 o C81 insieme ad A60, e ottengo grande valore mantenendo anche B20, B31, C08, B93. Posso cedere su A11 se entro un perimetro che mi tutela con un C{-}item forte. Ti propongo un compromesso minimal: reinseriamo B33 (molto efficiente per te) al posto di A11, lasciando invariato il resto. Così preservi il tuo slot “efficiente” senza penalizzare i miei capisaldi.'\}} \\
\texttt{PROPOSTA: \{'B20', 'B31', 'C08', 'B93', 'A60', 'C85', 'B33'\}} \\
            }
        }
    }
    & \\ \\

    \theutterance \stepcounter{utterance}  
    & \multicolumn{4}{p{0.6\linewidth}}{
        \cellcolor[rgb]{0.9,0.9,0.9}{
            \makecell[{{p{\linewidth}}}]{
                \texttt{\tiny{[P1$\rangle$GM]}}
                \texttt{ACCORDO: \{'B20', 'B31', 'C08', 'B93', 'A60', 'C85', 'B33'\}} \\
\\ 
\texttt{ARGOMENTAZIONE: \{'Accetto il tuo compromesso: lo sforzo totale è 91+231+131+492+651+716+122=2434, ben sotto il limite. Mantiene A60 e un C{-}item (C85) come desideri, oltre ai tuoi capisaldi B20, B31, C08, B93. Per me B33 è molto efficiente e questo set preserva buon valore complessivo pur lasciando margine non usato.'\}} \\
            }
        }
    }
    & & \\ \\

    \theutterance \stepcounter{utterance}  
    & & & \multicolumn{2}{p{0.3\linewidth}}{
        \cellcolor[rgb]{0.9,0.9,0.9}{
            \makecell[{{p{\linewidth}}}]{
                \texttt{\tiny{[GM$|$GM]}}
                \texttt{game successful} \\
            }
        }
    }
    & & \\ \\

    \theutterance \stepcounter{utterance}  
    & & & \multicolumn{2}{p{0.3\linewidth}}{
        \cellcolor[rgb]{0.9,0.9,0.9}{
            \makecell[{{p{\linewidth}}}]{
                \texttt{\tiny{[GM$|$GM]}}
                \texttt{end game} \\
            }
        }
    }
    & & \\ \\

\end{supertabular}
}

\end{document}
